\documentclass[11pt]{article}
\usepackage{mathunal-base}
\mathunalfuente{serif}
\providecommand{\nota}[1]{\par\smallskip\noindent{\small\itshape\color{unalblue}#1}\par\smallskip}
\newlist{opciones}{enumerate}{1}
\setlist[opciones]{label=(\roman*),itemsep=1pt,topsep=2pt,leftmargin=1.8em,font=\normalfont}
\newcommand{\rp}{\underline{\hspace{3cm}}}

\renewcommand{\examtitulo}{Segundo Examen Parcial de Ecuaciones Diferenciales}
\renewcommand{\examfecha}{Mayo del 2022}

\begin{document}

\examheader

\vspace{4pt}
\noindent\textit{Duración: 1 hora y 45 minutos. El examen consta de tres preguntas en dos hojas impresas por ambos lados; antes de empezar verifique que su examen esté completo y que sus datos sean correctos. No desengrape su examen ni añada otras grapas o su examen será anulado. En las preguntas con procedimiento justifique sus respuestas en los espacios asignados. En los puntos 1 y 2 debe justificar paso a paso su procedimiento; en el punto 3 sólo se calificará la respuesta. No está permitido hacer preguntas, usar calculadoras, sacar hojas en blanco ni ningún tipo de apuntes durante el examen. Por favor verifique que su celular esté apagado.}

\vspace{6pt}
\noindent\textbf{1. (25\%)} Halle la solución general de la ecuación diferencial
\[
y''-8y'+16y=x^{-1}e^{4x}
\]
en el intervalo $(0,\infty)$.

\vspace{7cm}

\newpage

\noindent\textbf{2. (20\%)} Una masa que pesa 24 lb estira un resorte $\frac{1}{12}$ pie. Suponga que el sistema oscila en un medio con una constante de amortiguamiento de $\beta$ lb$\cdot$s/pie. La masa se pone en movimiento desde un punto que está $\frac{2}{5}$ pie abajo de su posición de equilibrio, con una velocidad descendente de 7 pies/s. Suponga que sobre el sistema no actúan fuerzas externas.

\begin{opciones}
\item[a.] \textbf{(15\%)} Plantee el problema de valor inicial que satisface la posición $x(t)$ de la masa en el instante de tiempo $t$.
\item[b.] \textbf{(5\%)} Encuentre el valor de $\beta$ para el cual el sistema está críticamente amortiguado.
\end{opciones}

\nota{NOTA: No es necesario que resuelva este PVI, ni tal solución será tenida en cuenta.}

\vspace{5cm}

\newpage

\noindent\textbf{3. (55\%)} En los literales a), b), c) y d) complete según corresponda. Sólo se calificará la respuesta, la cual debe estar simplificada. No es necesario que escriba el procedimiento.

\vspace{0.25cm}
\noindent\textbf{a) (12\%)} Una forma adecuada de solución particular $y_p$ para la ecuación diferencial
\[
y^{(4)}+8y'''+16y''=7+8e^{-4t}+e^{-4t}\sin(2t), \qquad t\in\mathbb{R},
\]
es $y_p(t)=$~\rp\ (No halle las constantes).

\vspace{0.3cm}
\noindent\textbf{b) (13\%)} La solución general de la ecuación diferencial homogénea de Cauchy-Euler
\[
x^2y''+3xy'+5y=0, \qquad x>0,
\]
es
\[
y=\rp
\]

\vspace{0.3cm}
\noindent\textbf{c) (10\%)} Si $y_1$ y $y_2$ son soluciones L.I. en $(0,\infty)$ de la ecuación diferencial
\[
y''-\frac{3}{x}y'+xe^xy=0
\]
y si el Wronskiano de $y_1$ y $y_2$ evaluado en 2 es $W(y_1,y_2)(2)=1$, entonces
\[
W(y_1,y_2)(7)=\rp
\]

\vspace{0.3cm}
\noindent\textbf{d) (20\%)} La posición de un cuerpo, que está sujeto a un resorte, está dada por
\[
x(t)=-e^{-t}\cos(4t)+\sqrt3\,e^{-t}\sin(4t)
\]
pies, para cada instante $t\geq 0$ (medido en segundos). En los siguientes numerales complete según corresponda.

\begin{opciones}
\item \textbf{(12\%)} Al escribir la posición de la masa en la forma $x(t)=Ae^{-t}\sin(4t+\Phi)$, los valores de $A$ y $\Phi$ son respectivamente
\[
A=\rp\ \text{pies}, \qquad \Phi=\rp\ \text{radianes}.
\]
\item \textbf{(8\%)} El instante de tiempo $t$ en que la masa pasa por posición de equilibrio dirigiéndose hacia abajo por segunda vez es
\[
t=\rp\ \text{segundos}.
\]
\end{opciones}

\vfill
\rule{\linewidth}{0.4pt}\\[1pt]
{\scriptsize\textit{Transcripción digital --- MathUNAL}\hfill\textcolor{unalblue}{\textbf{1}}}

\end{document}
